\chapter*{Abstract}
\addcontentsline{toc}{chapter}{Abstract}

No finite substrate can be guaranteed sufficient for arbitrary future
continuation across an open task space: this is the informal thesis this
monograph makes precise and, in one central case, proves. Distinctions
survive transformation only through renderings that discard some of what
they carry; whether that discard is benign or corrupting depends on
whether it is declared, not merely on its size. When renderings compose
into a chain --- a substrate chain, or budget relay --- discard
accumulates, and the Composition Failure Theorem shows that a sequence of
locally rational transformations, each optimizing for its own task, can
jointly produce a substrate insufficient for a later task none of them
anticipated, without any single stage behaving negligently. The theorem
is proved against an explicit numerical witness and shown to be
independent of which local discard procedure a stage uses. A parallel
result identifies the opposite case: a chain sharing one task throughout,
adequately funded at every stage, cannot exhibit this failure, which is
part of why formal proofs behave so differently from informal
explanation under the same transmission process. The traditional
distinction between preserving an artifact and reconstructing it from
traces is replaced by a single continuous fidelity measure, with
preservation recovered only as its high-fidelity limit; a companion
chapter states plainly, rather than assumes away, the one dependency this
correction cannot settle on its own --- whether embodied perception ever
reaches that limit at all. Eight worked-domain chapters --- chess,
painting, maps, proofs, contracts, science, memory, and media pipelines
--- test the framework against cases chosen because they generated it,
not because they illustrate it after the fact. A late chapter situates this monograph
against the wider Distinguishability and Admissibility programs it draws
on, verifying prior claims of informal inheritance against primary
source text where that text was available, transcribing formulas
directly where only the author could supply them, and retracting one
citation to a text that does not appear to exist. Three findings survive
that audit intact: the framework's oldest working distinction already
had a name in a program it had not engaged; its treatment of
task-relative sufficiency was silently assuming an evaluability condition
a related program makes explicit; and its central theorem occupies a
setting a sibling theorem, proved for a single observer at one moment,
does not cover and cannot be reduced to.
